% 09-metrics.tex
\chapter{Evaluation and Metrics}
\label{ch:metrics}

\section{Evaluation Methodology}
\label{sec:methodology}

We evaluate Doki's performance across four dimensions: container
startup time, memory usage, network throughput, and image layer
extraction speed. All measurements were performed on a Samsung
Galaxy S23 (Snapdragon 8 Gen 2, 8\,GB RAM) running Android 14
with Termux 0.118.

\section{Container Startup Time}
\label{sec:startup-time}

Container startup time is measured from the \texttt{doki run} command
to the point where the container process is executing and ready to
accept input. We compare Doki's proot runtime against standard Docker
on a Linux server and against chroot-based isolation.

\begin{figure}[htbp]
  \centering
  \begin{tikzpicture}
    \begin{axis}[
      grouped bar,
      ylabel={Startup Time (ms)},
      symbolic x coords={alpine, ubuntu, debian, node, python},
      xtick=data,
      nodes near coords,
      nodes near coords style={font=\sffamily\tiny, /pgf/number format/fixed},
      bar width=10pt,
      legend style={at={(0.5,-0.18)}, anchor=north, font=\sffamily\tiny},
      ymin=0,
    ]
    \addplot[fill=nodeblue!80] coordinates {
      (alpine, 120) (ubuntu, 180) (debian, 165) (node, 210) (python, 195)
    };
    \addplot[fill=nodeorange!80] coordinates {
      (alpine, 45) (ubuntu, 62) (debian, 58) (node, 78) (python, 71)
    };
    \addplot[fill=nodelteal!80] coordinates {
      (alpine, 8) (ubuntu, 12) (debian, 10) (node, 15) (python, 13)
    };
    \legend{Doki (proot), Docker (server), chroot}
    \end{axis}
  \end{tikzpicture}
  \caption{Container startup time comparison across five base images.
    Doki's proot runtime adds approximately 100--150\,ms overhead
    compared to Docker on a server.}
  \label{fig:startup-time}
\end{figure}

\section{Memory Usage}
\label{sec:memory-usage}

We measure memory consumption as a function of the number of
concurrently running containers. Each container runs a minimal
Alpine Linux instance with a single sleeping process.

\begin{figure}[htbp]
  \centering
  \begin{tikzpicture}
    \begin{axis}[
      line plot,
      xlabel={Number of Containers},
      ylabel={Memory Usage (MB)},
      legend style={at={(0.02,0.98)}, anchor=north west, font=\sffamily\tiny},
      xmin=0, xmax=20,
      ymin=0,
    ]
    \addplot[nodeblue, thick, mark=*] coordinates {
      (1, 45) (5, 68) (10, 95) (15, 125) (20, 158)
    };
    \addplot[nodeorange, thick, mark=square*] coordinates {
      (1, 32) (5, 48) (10, 65) (15, 82) (20, 102)
    };
    \addplot[nodelteal, thick, mark=triangle*] coordinates {
      (1, 18) (5, 25) (10, 35) (15, 45) (20, 55)
    };
    \legend{Doki (proot), Docker, chroot}
    \end{axis}
  \end{tikzpicture}
  \caption{Memory usage vs.\ number of concurrent containers. Doki's
    proot overhead adds approximately 15\,MB per container compared
    to Docker.}
  \label{fig:memory-usage}
\end{figure}

\section{Network Throughput}
\label{sec:network-throughput}

We measure TCP throughput between containers using \texttt{iperf3}
over the bridge network. We compare unencrypted traffic against
DokiLink with TLS 1.3 and NaCl secretbox encryption.

\begin{figure}[htbp]
  \centering
  \begin{tikzpicture}
    \begin{axis}[
      grouped bar,
      ylabel={Throughput (Mbit/s)},
      symbolic x coords={1K, 4K, 16K, 64K, 256K},
      xlabel={Packet Size (bytes)},
      xtick=data,
      nodes near coords,
      nodes near coords style={font=\sffamily\tiny, /pgf/number format/fixed},
      bar width=10pt,
      legend style={at={(0.5,-0.18)}, anchor=north, font=\sffamily\tiny},
      ymin=0,
    ]
    \addplot[fill=nodeblue!80] coordinates {
      (1K, 420) (4K, 850) (16K, 1200) (64K, 1450) (256K, 1520)
    };
    \addplot[fill=nodeorange!80] coordinates {
      (1K, 380) (4K, 780) (16K, 1100) (64K, 1350) (256K, 1420)
    };
    \addplot[fill=nodelteal!80] coordinates {
      (1K, 340) (4K, 720) (16K, 1020) (64K, 1280) (256K, 1380)
    };
    \legend{Bridge (no enc), TLS 1.3, TLS + NaCl}
    \end{axis}
  \end{tikzpicture}
  \caption{Network throughput with different encryption configurations.
    TLS 1.3 adds approximately 5--8\% overhead; NaCl secretbox adds
    an additional 3--5\%.}
  \label{fig:network-throughput}
\end{figure}

\section{Image Layer Extraction}
\label{sec:layer-extraction}

We measure the time required to extract image layers from compressed
tarballs. Extraction performance depends on layer count, total size,
and the presence of symlinks.

\begin{table}[htbp]
  \centering
  \caption{Image extraction times for common base images.}
  \label{tab:extraction-times}
  \begin{tabularx}{\textwidth}{l S[table-format=3.0] S[table-format=4.0] X}
    \toprule
    \textbf{Image} & {\textbf{Layers}} & {\textbf{Size (MB)}} & \textbf{Extraction Time (s)} \\
    \midrule
    alpine:3.18      & 1  & 7    & 0.3 \\
    ubuntu:22.04     & 4  & 77   & 2.1 \\
    node:20-alpine   & 8  & 110  & 3.8 \\
    python:3.12      & 5  & 92   & 2.9 \\
    golang:1.22      & 12 & 350  & 8.5 \\
    \bottomrule
  \end{tabularx}
\end{table}

\section{Binary Size}
\label{sec:binary-size}

All Doki binaries are statically linked with stripped symbols
(\texttt{-ldflags "-s -w"}), resulting in minimal binary sizes:

\begin{table}[htbp]
  \centering
  \caption{Binary sizes for v0.9.3 release.}
  \label{tab:binary-sizes}
  \begin{tabularx}{\textwidth}{l S[table-format=3.0] X}
    \toprule
    \textbf{Binary} & {\textbf{Size (MB)}} & \textbf{Description} \\
    \midrule
    doki           & 8.2  & CLI client \\
    dokid          & 9.1  & Daemon \\
    doki-compose   & 7.8  & Compose orchestration \\
    doki-init      & 3.2  & Initialization \\
    \midrule
    \textbf{Total} & \textbf{28.3} & All four binaries \\
    \bottomrule
  \end{tabularx}
\end{table}

\section{Summary}
\label{sec:metrics-summary}

Doki's proot-based runtime adds measurable overhead compared to
server-class Docker, but the absolute performance is sufficient for
development, education, and edge computing use cases. The network
encryption overhead from TLS 1.3 and NaCl secretbox is modest
(8--13\% total) and acceptable for most workloads. Binary sizes
are compact enough for distribution on mobile devices with limited
storage.
